%% \exannot: the annotation column.  Four things fail silently here, and
%% each has its own examples.
%%
%%   * The column is ONE column across NESTING DEPTHS.  \linewidth has lost
%%     the indentation by the time an example body is set, so a column
%%     measured from it drifts one indent step per level -- 247 / 271 / 295pt
%%     for the three levels here, which looks entirely deliberate on the
%%     page.  ANNTOP / ANNSUB / ANNDEEP.
%%   * The column is ONE column across TEXT LENGTHS.  This is where \jambox
%%     gives way: it sets its box at its natural width behind glue that
%%     shrinks to nothing, so an example that reaches the column pushes the
%%     annotation right rather than wrapping, with no overfull warning to
%%     say so.  ANNLONG does reach it: its annotation must be on a line of
%%     its own AND still in the column.
%%   * \ExAnnotSep is a LEAST gap, and is rigid so as to be one.  ANNTIGHT
%%     is built to end .4em short of the column -- \settowidth, so it is
%%     .4em short under every engine and not by luck of the metrics.  Its
%%     annotation must wrap.  With shrink in that gap it would not: it
%%     would stay on the line, still in the column (the fixed-width box
%%     sees to that), but nearer the example text than \ExAnnotSep, up to
%%     touching it.  So this is the one example that tells a rigid gap from
%%     a rubber one, and nothing else in the suite does.
%%   * A gloss annotation belongs on the OBJECT line, level with the
%%     language it labels -- not under the free translation, and not as a
%%     column of the grid, which would push the gloss tier beneath it.
%%     ANNGLOSS must share a line with the object tier and not with
%%     ANNGLOSSTIER, which is the gloss directly under it.
%%   * The filler is a fill and not a fil.  A gloss is set \raggedright, so
%%     \rightskip is already 0pt plus 1fil; an \hfil filler would share the
%%     slack with it instead of taking it, and the gloss annotation would
%%     come to rest partway to the column.  That mutation is invisible in
%%     every example above -- \rightskip is zero in all of them -- and is
%%     caught by ANNGLOSS alone.  \ExRaggedRight ([langsci]) is the same
%%     path with the same filler.
%%
%%   * \exannot ends its paragraph, and a blank line ends a dot-syntax
%%     example: ANNHEAD pins that the one is not the other.
%%
%% Spoken forms and the structure tree are exannot-ua.tex; this case is
%% untagged and about the geometry alone.
\input{_preamble}
\setlength{\ExAnnotColumn}{9cm}
\begin{document}
\ex. ANNTOP a top-level example\exannot{[XP]}

%% An annotated HEAD with sub-examples under it.  \exannot ends the
%% paragraph -- it has to, or \parfillskip could not hold the last line to
%% the column -- and a blank line is what ends a dot-syntax example, so the
%% \par is one token away from closing the example instead of the line.
%% If it ever does, the \a. below has no example to attach itself to and
%% the case stops compiling; the numbering is asserted too, since an
%% example closed early would take a number of its own.
\ex. ANNHEAD an annotated head\exannot{[HP]}
\a. ANNUNDER a sub-example under it\exannot{[IP]}
\z.

\ex.
\a. ANNSUB a sub-example\exannot{[CP]}
\b. ANNSIB and its sibling\exannot{[TP]}
\a. ANNDEEP one level deeper\exannot{[DP]}
\b. ANNLONG a sub-example long enough to reach the column and then some\exannot{[VP]}
\z.

\newlength{\anntightw}
\begingroup
\settowidth{\anntightw}{ANNTIGHT a snug example}
\setlength{\ExAnnotColumn}{\dimexpr\Exindent+\Exlabelwidth+\Exlabelsep
                                  +\anntightw+.4em\relax}
\ex. ANNTIGHT a snug example\exannot{[QP]}
\endgroup

\ex. \gll ANNGLOSS Pierre est fatigu\'e \exannot{[WP]}\\
          ANNGLOSSTIER Pierre is tired\\
     \glt `Pierre is tired'

%% \exannot GLUED to the last object word, with no space in front of it.
%% Both spellings of a glossed sub-example, because they reach \gll by
%% different routes: \b. followed by \gll, and the \bg. shorthand, which is
%% \b. plus the gloss head.
%%
%% Two things fail here and only one of them is loud.
%%
%% Loud: without the space "libro\exannot{...}" is ONE word of tier 1 and
%% its head token is "l", so a rule that looks only at the head of the word
%% never sees the annotation, \exannot runs where it stands, and the run
%% stops on an error about where annotations have to go -- which is not
%% what was wrong with it.  Reported from a beamer deck.
%%
%% Silent, and worse: the word is split at the token, and a DELIMITED
%% argument that is exactly one brace group has its braces removed by TeX.
%% So the tail came back as "(italien)" rather than "{(italien)}", the
%% annotation was rebuilt as "\exannot (italien)", its mandatory argument
%% was the single token "(", and the page carried "(" in the column with
%% "italien)" flung to the right margin.  Zero errors, zero warnings.
%% GLUEA and GLUEB are asserted to sit in the column WHOLE, which is the
%% only way to tell that apart from a correct one.
\ex.
\a. GLUEPLAIN le livre / mon livre / *le mon livre
\b. \gll GLUEA il mio libro\exannot{(italien)}\\
         GLUEATIER le mon livre\\
    \glt `mon livre'
\z.

\ex.
\a. GLUEPLAINB le livre / mon livre / *le mon livre
\bg. GLUEB il mio libro\exannot{(italien)}\\
     GLUEBTIER le mon livre\\
\glt `mon livre'
\z.

%% An annotation whose ARGUMENT CONTAINS SPACES, which is what prose does:
%% "(cf. \emph{le mien livre})" has three of them.  \gll splits the object
%% line on spaces, so an annotation still sitting in that line when the
%% splitter runs is not one thing by the time anyone looks: it becomes the
%% words "libro\exannot{(cf.", "\emph{le", "mien", "livre})", each of them a
%% COLUMN of the grid, with gloss cells appearing underneath the pieces.
%%
%% What that looks like on the page is not "my annotation came apart" but
%% "the gloss has a line break in it", because the extra columns stop
%% fitting and the grid wraps -- which is how it was reported, from a class
%% deck.  Nothing errors.  SPCA is asserted to reach the page as ONE word at
%% the column, which is the only thing that tells it from a gloss that
%% happens to have more words in it.
\ex.
\a. SPCPLAIN le livre / mon livre / *le mon livre
\b. \gll SPCOBJ il mio libro\exannot{(cf. \emph{le mien livre})}\\
         SPCTIER le mon livre\\
    \glt `mon livre'
\z.

%% A gloss with NO annotation, standing right after one that has it.  The
%% lifted annotation is stashed in a token list, and a token list that is
%% not cleared between glosses is re-emitted by the next one: an annotation
%% appears against an example that never asked for it, in the right column
%% and the right font, looking entirely deliberate.  SPCNONE must carry
%% nothing.
%% In the SAME example as an annotated one, deliberately: the stash is a
%% local assignment, so a plain gloss in a LATER example is protected by
%% the group boundary between them and shows nothing.  Two sub-examples of
%% one example share a group, and that is where a stale stash surfaces.
\ex.
\a. \gll SPCKEEP il mio libro\exannot{(voir \emph{le mien livre})}\\
         SPCKEEPTIER le mon livre\\
\b. \gll SPCNONE il mio libro\\
         SPCNONETIER le mon livre\\
\z.

%% An annotation too wide to share the object line must move ITSELF down.
%% It must not make room by breaking the GRID -- which is what it did, and
%% what was reported from a class deck:
%%
%%     il  mio                            il  mio  libro molto lungo
%%     le  mon               instead of   le  mon  livre tres long
%%     libro molto lungo                      (a ; cf. a. fr. ...)
%%     livre tres long
%%
%% A gloss is set \raggedright, so every way of breaking it has badness
%% zero and the penalty is the only thing telling one from another.  The
%% glue between two columns used to be a FREE breakpoint while \exannot
%% offers its own at \penalty50, so TeX took the free one and split the
%% object line between two words.  \lx@gl@colpenalty is what makes the grid
%% dearer than the annotation.
%%
%% Whether it triggers is a WINDOW and not a threshold, which is why it went
%% unnoticed for so long: the annotation has to be too wide for the room the
%% column leaves AND not so wide that breaking the grid stops making room
%% for it either.  The width of that window is the width of the LAST gloss
%% column, so the column here is deliberately a long braced one -- with a
%% short last column the window is a few points across and which side of it
%% a given annotation falls on is luck.  An earlier version of this case had
%% exactly that, sat outside the window, and passed with the penalty
%% removed.  Two annotation widths, so one lucky width cannot carry it.
\ex.
\a. BRKA le livre
\b. \gll BRKB il mio {libro molto lungo veramente} \exannot{(a ; cf. a. fr. le mien livre)}\\
         gl le {mon livre vraiment tres long}\\
\z.

\ex.
\a. BRKAC le livre
\b. \gll BRKC il mio {libro molto lungo veramente} \exannot{(b ; cf. a. fr. le mien livre, bien plus long que cela)}\\
         gl le {mon livre vraiment tres long}\\
\z.

%% ... and a braced multi-word object item immediately before a glued
%% annotation.  The braces there are the word splitter's, and losing them
%% at the split would turn one object word into two and pull the gloss out
%% of line under it.
\ex. \gll GLUEC ein {ganz kurzes} GOBJ\exannot{(gebraucht)}\\
          GLUECT a {very short} GGLOSS\\
\end{document}
